\documentclass[10pt]{article}
\usepackage[margin=1in]{geometry}
\usepackage{booktabs,amsmath,microtype,enumitem}
\setlist{itemsep=2pt,topsep=3pt}
\title{Finding Object States in Cooking Videos Without a State List}
\author{ChefOST (experiment, not published results)}
\date{October 2026}
\begin{document}
\maketitle

\section{Method}
\textbf{Goal.} Given a cooking video and the name of an object in it (e.g.\ ``bowl''), find the states the
object goes through and when each one happens. The number of states is not fixed. Each state comes out as an
embedding vector, plus a short text description if wanted. We never give the pipeline a list of possible
states, and we train nothing. MOSCATO labels are used only to score the output.

The pipeline has three steps.

\textbf{Step 1: turn the video into a track of embeddings.} We take one frame per second and embed it with
SigLIP (SO400M/14-384), an image--text model. This gives one vector $x_t$ per second. We tried two inputs:
\begin{itemize}
\item \emph{Whole frame:} embed the full image.
\item \emph{Object crop:} find the object with the OWLv2 detector, prompted only with its name, and embed a
crop around the best box (25\% margin). The detector's confidence tells us how visible the object is. We turn
it into a weight $w_t$ between 0 and 1 by rescaling between that track's 20th and 90th percentile scores.
(For whole frames, $w_t=1$.)
\end{itemize}

\textbf{Step 2: split the track into states.} An object's state shows up as a stretch of time where its
embedding stays roughly the same.
\begin{enumerate}
\item \emph{Smooth.} Average each frame with its neighbours (Gaussian, $\sigma=3$\,s), weighting by
visibility. Frames where the object is hidden take their value from nearby frames where it is seen.
\item \emph{Cut into segments.} Find the best split of the track into segments of at least 5\,s, where each
segment is summarised by its average. We minimise how far frames are from their segment average, plus a
fixed cost $\beta$ for each extra segment:
\[ \textstyle\sum_{\text{segments } j}\;\sum_{t\in j} w_t\,\|x_t-\mu_j\|^2 \;+\; \beta\,\hat\sigma^2 \cdot (\text{number of segments}). \]
Here $\hat\sigma^2$ measures the track's frame-to-frame noise (half the mean squared change between
consecutive frames), so $\beta$ works the same for noisy and quiet tracks. A dynamic program finds the exact
best split quickly.
\item \emph{Merge.} Segments whose averages are very similar (cosine distance below $\tau$) become one state,
so a state can come back later (e.g.\ empty $\to$ full $\to$ empty).
\end{enumerate}
Each state is represented by the weighted average $c_k$ of its frames.

\textbf{Step 3 (optional): describe each state with a VLM.} For each state we pick two representative frames:
one from each half of the state's time span, choosing the frame that is most visible and closest to the
state's average. We ask Qwen2.5-VL-3B for a phrase of at most 8 words describing the object's state. The
prompt mentions only generic aspects (shape, cut, mixed, cooked, full, clean, location) and no MOSCATO words.
Each phrase $d$ is embedded with SigLIP's text encoder as ``a photo of \{object\}, $d$.''

\textbf{What each frame outputs.} Every frame outputs the vector of the state it belongs to. We compare three
choices:
\begin{itemize}
\item \textsc{discover}: the state's image average $c_k$, minus the average image embedding $\mu$ of all
frames in the split. This corrects for image and text embeddings sitting in different regions of the shared
space (the ``modality gap''). No labels are needed to compute $\mu$.
\item \textsc{vlm}: the average text embedding of the state's descriptions.
\item \textsc{fused}: the sum of the two (each normalised to length 1).
\end{itemize}

\section{Experiments}
\textbf{Data.} 30 CMU-MMAC cooking videos with MOSCATO ground truth (GT), 267 video--object pairs, one frame
per second. We split the videos into:
\begin{itemize}
\item \emph{Dev}: 2 videos per recipe (10 videos, 79 pairs), used to choose settings ($\beta=500$,
$\tau=0.05$, smoothing, modality-gap correction).
\item \emph{Test}: the other 20 videos (188 pairs, 23{,}074 scored frames), used for the reported numbers.
\end{itemize}
Confidence intervals (95\%) come from resampling videos (paired bootstrap).

\textbf{How we score state names.} To compare vectors with GT labels, we embed each possible state of the
object as text (``a photo of \{object\}, \{state\}.'') with the same SigLIP text encoder. The list of possible
states comes from MOSCATO pseudo-labels of videos outside the GT set, and is used only for scoring. It
contains at least one correct state on 99.6\% of frames. On each frame where the object has a GT state, we
rank the possible states by similarity to the output vector:
\begin{itemize}
\item \emph{Top-1}: the highest-ranked state is correct.
\item \emph{MRR}: 1 divided by the rank of the first correct state (1st $=1$, 2nd $=0.5$, \dots).
\item \emph{P@2 / R@2}: treat the top 2 states as the prediction and compute the README precision and recall.
\end{itemize}
``Chance'' is the exact expected score if the states were ranked at random.

\textbf{How we score the discovered states.} We compare the discovered segments with GT segments (stretches
where the GT state set stays the same):
\begin{itemize}
\item \emph{MoF}: fraction of frames correct after matching each discovered state to one GT state.
\item \emph{NMI, ARI}: standard measures of how well two groupings agree (0 = none, 1 = perfect). Computed
only on pairs whose GT state changes at least once.
\item \emph{Change F1}: whether discovered state changes fall within 5\,s of real ones.
\item \emph{\#states}: average number of states found per pair.
\end{itemize}

\begin{table}[h]\centering\small
\begin{tabular}{llcccccc}
\toprule
 & & \multicolumn{4}{c}{Test (20 videos)} & \multicolumn{2}{c}{Dev (10)}\\
Input & Output & Top-1 & MRR & P@2 & R@2 & Top-1 & MRR\\
\midrule
--- & chance & .171 & .361 & .171 & .157 & .164 & .359\\
InternVideo frame & per frame & .319 & .495 & .279 & .249 & .259 & .435\\
SigLIP frame & per frame & .353 & .521 & .298 & .264 & .278 & .471\\
SigLIP frame & one state per pair & .337 & .507 & .293 & .256 & .305 & .463\\
SigLIP frame & \textsc{discover} & .369 & .539 & .316 & .283 & .293 & .494\\
SigLIP crop & per frame & .281 & .470 & .257 & .234 & .250 & .448\\
SigLIP crop & \textsc{discover} & .301 & .493 & .277 & .251 & .314 & .504\\
SigLIP frame & \textsc{vlm}, evenly spaced frames & .362 & .513 & .299 & .275 & .336 & .506\\
SigLIP frame & \textsc{vlm}, frames picked per state & .375 & .523 & .311 & .288 & .315 & .489\\
SigLIP crop & \textsc{fused} & .394 & .553 & .317 & .302 & .255 & .460\\
SigLIP frame & \textsc{fused} & \textbf{.438} & \textbf{.580} & \textbf{.346} & \textbf{.312} & .313 & .495\\
\bottomrule
\end{tabular}
\caption{How well each output names the object's state. ``Evenly spaced frames'' uses the same number of VLM
calls as ``frames picked per state'', but spreads them evenly over time instead.}
\end{table}

\textbf{Results: naming states (Table 1, test).}
\begin{itemize}
\item \textsc{fused} is best: top-1 .438 vs .171 chance, MRR .580 vs .361.
\item It beats one embedding per frame by $+0.086$ top-1 (CI $+0.046$ to $+0.123$) and $+0.058$ MRR
($+0.025$ to $+0.093$).
\item It also beats each of its parts: \textsc{discover} by $+0.069$ ($+0.017$ to $+0.119$) and \textsc{vlm}
by $+0.064$ ($+0.027$ to $+0.103$).
\item Grouping frames into states helps a little on its own: $+0.018$ MRR ($+0.002$ to $+0.034$).
\item Picking VLM frames per state is not clearly better than evenly spaced frames ($+0.012$, $-0.015$ to
$+0.038$).
\item The VLM is called only about 7 times per pair (about 2\% of frames), because it describes states, not
frames.
\end{itemize}

\begin{table}[h]\centering\small
\begin{tabular}{llccccc}
\toprule
Input & Segments & MoF & NMI & ARI & Change F1 & \#states (GT 2.90)\\
\midrule
--- & one state & \textbf{.735} & .000 & .000 & .000 & 1.00\\
SigLIP crop & discovered & .729 & .399 & .314 & \textbf{.331} & \textbf{3.02}\\
SigLIP crop & evenly spaced, same number & .647 & .425 & .297 & .247 & 3.66\\
SigLIP frame & discovered & .694 & .408 & .317 & .324 & 3.67\\
SigLIP frame & evenly spaced, same number & .538 & .449 & .266 & .276 & 5.88\\
InternVideo frame & discovered & .667 & \textbf{.449} & \textbf{.326} & .321 & 4.72\\
\bottomrule
\end{tabular}
\caption{How well the discovered states match GT segments (test). ``Evenly spaced'' cuts the track into as
many equal pieces as discovery found before merging, so it tests whether discovery gets the timing right.}
\end{table}

\textbf{Results: finding states (Table 2, test).}
\begin{itemize}
\item With object crops, discovery finds 3.02 states per pair, close to GT's 2.90.
\item It places state changes better than evenly spaced cuts (Change F1 .331 vs .247).
\item One constant state per pair still scores slightly higher on MoF (.735 vs .729), because most MOSCATO
objects barely change state.
\end{itemize}

\textbf{Caveats.}
\begin{itemize}
\item \emph{Dev and test disagree.} On dev, \textsc{fused} gains only $+0.035$ top-1 ($-0.015$ to $+0.086$,
not significant), and crops do better than whole frames. On test, whole frames do better. Dev has only 10
videos.
\item \emph{The best row was chosen after seeing test.} The whole-frame \textsc{fused} run was added after the
crop results on test. The crop version, planned in advance, also beats one embedding per frame on test
($+0.042$, $-0.001$ to $+0.086$).
\item \emph{Shared encoder.} The pipeline and the scorer both use SigLIP's text encoder.
\item \emph{Whole frames can track the wrong thing.} They follow the whole scene, not just the object. In
video S14 Eggs, parts of the \emph{bowl}'s track are described as ``pan with egg mixture cooking''.
\end{itemize}

\section{Conclusion}
Object states can be found without any state list by splitting an object's embedding track over time. The
number of states found is close to GT (3.0 vs 2.9 per pair), and the timing of changes beats evenly spaced
cuts. Describing each state once with a small VLM and adding that description to the visual average gives
the best results: top-1 .438 (chance .171) and MRR .580 (chance .361), with about 7 VLM calls per pair.

Because dev and test disagree and the best setup was chosen after seeing test, the size of this gain still
needs to be confirmed on more videos (EgoPER, EGTEA). The next steps are better object detection, so that
crops beat whole frames, and outputting several states per frame, since GT usually lists 2 or more.

\end{document}
